% !PRINTED 175
\weberpage{175}

\begin{center}
§. 54.

\gesperrt{The Tschirnhausen Transformation of the Equation of the Fifth Degree.}
\end{center}

If one wishes to apply the Tschirnhausen transformation to the equation of the fifth degree
\weq{1}{x^5 + a_1 x^4 + a_2 x^3 + a_3 x^2 + a_4 x + a_5 = 0}
one has to set
\weq{2}{y = \alpha_0 + \alpha_1 x + \alpha_2 x^2 + \alpha_3 x^3 + \alpha_4 x^4}
and obtains the transformed equation
\weq{3}{y^5 + b_1 y^4 + b_2 y^3 + b_3 y^2 + b_4 y + b_5 = 0.}

The most obvious thought would now be to determine the ratios of the
five unknowns $\alpha$ in such a way that
\weq{4}{b_1 = 0, \quad b_2 = 0, \quad b_3 = 0, \quad b_4 = 0}
would hold, which would reduce (3) to a pure equation; and
this was no doubt also the goal that \gesperrt{Tschirnhausen} had in
view. But of the four equations (4) the first is linear,
the second quadratic, the third of the third and the fourth of the
fourth degree. The final equation that one can derive from them
will therefore (by \S.\,50) be of degree $2.3.4 = 24$, and hence
no use can be drawn from it for the solution of the equation of the
fifth degree. One therefore seeks at first to satisfy only the
two equations
\weq{5}{b_1 = 0, \quad b_2 = 0,}
whereby the equation (3) passes into a form which, following
\gesperrt{F. Klein} (Lectures on the Icosahedron), is called a \gesperrt{principal equation} of the fifth degree.

To satisfy the two equations (5) we have at our
disposal the four ratios $\alpha_0 : \alpha_1 : \alpha_2 : \alpha_3 : \alpha_4$, and we
can still use these quantities in manifold ways to simplify the equation of the fifth degree.

We shall return to this subject in more detail in a later
section, and confine ourselves here to indicating the
aims in general terms. For the sake of vividness
we make use of a geometrical dress, which, incidentally,
is by no means essential to the understanding of the algebraic theory
as it is to be given later.
